% !TEX program = xelatex
% Technical Report 2.0, generated from the reviewed local artifacts.
\documentclass[11pt,a4paper]{article}
\usepackage[a4paper,margin=20mm]{geometry}
\usepackage{fontspec}
\usepackage{kotex}
\usepackage{booktabs,longtable,tabularx,array}
\usepackage[table]{xcolor}
\usepackage{enumitem}
\usepackage{amsmath}
\usepackage{fancyhdr}
\usepackage{graphicx}
\usepackage{float}
\usepackage{pgfplots}
\usepackage[hidelinks]{hyperref}
\pgfplotsset{compat=1.18}
\setmainfont{Noto Serif CJK KR}
\setsansfont{Noto Sans CJK KR}
\setmonofont{Noto Sans Mono CJK KR}
\definecolor{navy}{HTML}{173458}
\definecolor{blue}{HTML}{3867D6}
\definecolor{gold}{HTML}{C48A18}
\definecolor{orange}{HTML}{D97706}
\definecolor{ink}{HTML}{273142}
\definecolor{muted}{HTML}{667085}
\definecolor{grid}{HTML}{D0D5DD}
\definecolor{light}{HTML}{F5F7FA}
\newcommand{\code}[1]{\texttt{#1}}
\newcolumntype{Y}{>{\raggedright\arraybackslash}X}
\pagestyle{fancy}
\setlength{\headheight}{14pt}
\fancyhf{}
\lhead{RAI Risk Taxonomy}
\rhead{Technical Report 2.0}
\cfoot{\thepage}
\hypersetup{pdftitle={RAI Risk Taxonomy Technical Report 2.0},pdfauthor={RAI Risk Taxonomy Project}}

\title{\sffamily\bfseries RAI Risk Taxonomy\\[4pt]\Large Technical Report 2.0}
\author{1,711개 등록 ID와 1,660개 활성 L4의 분류 및 통계 검증}
\date{2026년 7월 25일}

\begin{document}
\maketitle

\begin{center}
\fcolorbox{navy}{light}{\parbox{0.92\linewidth}{
\textbf{현재 릴리스 후보.} 1,711개 등록 ID를 보존하고, 1,660개 활성 카드를
54개 의미 L3에 연결한다. 이 중 614개는 HOLD 검토 상태로 유지한다. HOLD는
별도의 의미 범주가 아니라 사람의 분류 검토를 위한 운영 오버레이이다.}}
\end{center}

\section{기술 요약}

본 보고서는 Physical AI의 기존 182개 분류를 정답 잠금 집합으로 유지하면서,
글로벌 L4 1,726개를 신규 50개 L3에 연결한 단계적 절차를 기술한다. 기존 글로벌
L1\textendash{}L3는 예측 특징으로 사용하지 않았고 영구 L4 ID와 교차매핑을 보존했다.

Stage 1은 엄격한 rule, semantic, anchor 합의와 두 전문 에이전트의 exact-L3 합의를
요구해 Physical 182개와 신규 제안 37개만 배치했다. Stage 2는 hierarchy-blind
top-1 후보를 적극 수용해 1,553개(89.9768\%)를 배치하고 173개를 보류했다.
Stage 3 실험은 173개를 모두 강제 배치했으나, 정책 검토 결과 이 중 hard hold
55개를 다시 분리했다. 최종 정책 뷰는 1,671개 분류와 55개 검토 보류로 구성된다.

\begin{figure}[H]
\centering
\begin{tikzpicture}
\begin{axis}[
  ybar, ymin=0, ymax=1800, width=0.92\linewidth, height=7cm,
  bar width=20pt, ylabel={카드 수},
  symbolic x coords={Stage 1,Stage 2,Stage 3 실험,최종 정책}, xtick=data,
  nodes near coords, nodes near coords align={vertical},
  axis line style={ink}, tick label style={font=\small},
  ymajorgrids, grid style={grid}, legend style={at={(0.5,-0.18)},anchor=north,legend columns=2},
]
\addplot[fill=blue,draw=navy] coordinates {(Stage 1,219) (Stage 2,1553) (Stage 3 실험,1726) (최종 정책,1671)};
\addplot[fill=gold,draw=orange] coordinates {(Stage 1,1507) (Stage 2,173) (Stage 3 실험,0) (최종 정책,55)};
\legend{L3 배치,보류}
\end{axis}
\end{tikzpicture}
\caption{단계별 L3 배치와 보류. 분모는 각 단계에서 동일한 1,726개다.}
\end{figure}

Stage 3의 100\%는 운영상 경로 완전성을 시험한 결과이고, 최종 정책의 96.8134\%는
명백한 경계 충돌을 별도 검토하도록 선택한 결과다. 두 수치를 정확도 또는 승인률로
해석해서는 안 된다.

\section{범위, 데이터와 식별 원칙}

분석 단위는 영구 ID \code{RAI4-0001}\textendash{}\code{RAI4-1726}의 L4 카드다. Physical AI
사이트의 182개 분류는 잠금 정답이고, 동일 이름이 여러 경로에 존재할 수 있으므로
라벨을 키로 사용하지 않는다. 상위 분류체계는 L0 1개, L1 3개, L2 6개, L3 50개로
구성된다. \code{RAI-HOLD}는 운영 버킷이며 L3 노드가 아니다.

\begin{table}[H]
\centering
\begin{tabular}{lrr}
\toprule
최종 정책 상태 & 건수 & 전체 비율 \\
\midrule
Physical 정답 잠금 & 182 & 10.5446\% \\
Stage 1 전문가 합의 제안 & 37 & 2.1437\% \\
Stage 2 알고리즘 제안 & 1,334 & 77.2885\% \\
Stage 3 잔여 강제 후보 유지 & 118 & 6.8366\% \\
Taxonomy Decision Hold & 55 & 3.1866\% \\
\midrule
합계 & 1,726 & 100.0000\% \\
\bottomrule
\end{tabular}
\caption{최종 정책 뷰의 상태 구성.}
\end{table}

\section{단계별 분류 절차와 기준}

\subsection{Stage 1: 정확 적합성 우선}

비-Physical 입력에는 계층 재진술을 제거한 L4 라벨, 직접 메커니즘 정의와 evidence
title만 사용했다. 자동 후보는 단일 rule eligibility, rule/semantic/composite top-1
일치, semantic score $\geq0.55$, composite score $\geq0.54$, 두 margin
$\geq0.035$, anchor 2/3 이상 합의를 모두 만족해야 했다. gap sentinel과 다중
메커니즘을 배제하고 두 전문 에이전트가 동일한 exact L3를 승인한 37개만 제안했다.

\subsection{Stage 2: 약 90\% 커버리지}

Stage 2는 보호 집합 219개를 변경하지 않고 다음 순위 점수를 사용했다.
\[
S_2=0.55q_{sem}+0.20q_{sm}+0.10q_{cm}+0.10q_{vote}+0.05q_{rule}-p_{gap}.
\]
여기서 의미 점수와 두 margin은 사전 고정 구간으로 0\textendash{}1 정규화했고, anchor vote는
3으로 나눴으며, rule support는 이진값, gap penalty는 0.05다. hard hold 55개를
먼저 보존하고 $S_2$ 하위 118개를 추가 보류해 총 173개(10.0232\%)를 남겼다.
나머지 1,334개는 hierarchy-blind top-1 L3에 배치했다.

\subsection{Stage 3와 최종 보류 정책}

Stage 3 실험은 잔여 173개를 기존 top-1에 강제 연결해 100\% 경로를 시험했다.
그러나 다음 다섯 hard-hold 사유의 55개는 자동 배치의 의미적 또는 정책적 위험이
크므로 최종 정책 뷰에서 \code{policy\_l3\_id=null}로 되돌리고
\code{RAI-HOLD}로 분리했다. 나머지 quota hold 118개는 강제 후보 상태와 사람
검토 의무를 유지한다.

\begin{figure}[H]
\centering
\begin{tikzpicture}
\begin{axis}[
  xbar, xmin=0, xmax=125, width=0.70\linewidth, height=7cm,
  xlabel={카드 수},
  symbolic y coords={절대 적합도 부족,다중 메커니즘,전문가 불일치,전문가 거절,Physical 범위 검토,적합도 reserve},
  ytick=data, nodes near coords, nodes near coords align={horizontal},
  bar width=12pt, xmajorgrids, grid style={grid}, axis line style={ink},
]
\addplot[fill=gold,draw=orange] coordinates {(2,절대 적합도 부족) (6,다중 메커니즘) (7,전문가 불일치) (17,전문가 거절) (23,Physical 범위 검토)};
\addplot[fill=blue,draw=navy] coordinates {(118,적합도 reserve)};
\end{axis}
\end{tikzpicture}
\caption{Stage 2 보류 173개의 사유. 금색 55개는 최종 \code{RAI-HOLD}, 파란색 118개는 Stage 3 강제 후보 유지 대상이다.}
\end{figure}

Physical 범위 검토 23개는 Physical AI로 확정된 카드가 아니다. robot, drone,
autonomous vehicle, embodied AI, cyber-physical, physical actuator 등의 표현이
탐지됐지만 잠긴 Physical 정답 182개와 동일하지 않아 경계 검토가 필요한 카드다.

\section{최종 분포 통계}

\subsection{L1과 L2 분포}

보류 55개를 제외한 1,671개 중 일반 AI는 1,140개, 에이전틱 AI는 349개,
피지컬 AI는 182개다. 분류된 카드 내 비율은 각각 68.2226\%, 20.8857\%,
10.8917\%다. 전체 1,726개의 나머지 3.1866\%는 보류다.

\begin{figure}[H]
\centering
\begin{tikzpicture}
\begin{axis}[
  xbar, xmin=0, xmax=1250, width=0.92\linewidth, height=5.7cm,
  xlabel={분류된 카드 수}, symbolic y coords={Physical AI,Agentic AI,General},
  ytick=data, nodes near coords, bar width=18pt, xmajorgrids, grid style={grid},
]
\addplot[fill=blue,draw=navy] coordinates {(182,Physical AI) (349,Agentic AI) (1140,General)};
\end{axis}
\end{tikzpicture}
\caption{최종 정책 뷰의 L1 분포. 분모는 L3 배치가 있는 1,671개다.}
\end{figure}

\begin{longtable}{llrr}
\toprule
L2 ID & L2 이름 & 카드 수 & 분류 1,671개 내 비율 \\
\midrule
\endhead
RAI2-G-INT & General Interaction & 708 & 42.3698\% \\
RAI2-G-SYS & General System & 432 & 25.8528\% \\
RAI2-A-SYS & Agentic System & 349 & 20.8857\% \\
RAI2-P-SYS & Physical System Safety & 91 & 5.4458\% \\
RAI2-P-INT & Physical Interaction Safety & 62 & 3.7104\% \\
RAI2-P-SOC & Physical Societal Safety & 29 & 1.7355\% \\
\bottomrule
\end{longtable}

\subsection{L3 집중도}

50개 L3 모두 최소 한 카드를 포함한다. 가장 큰 L3는 Anthropomorphism 242개로
분류된 1,671개의 14.4823\%다. 상위 5개는 700개(41.8911\%)이며 단일 L3가
20\%를 넘지 않는다. 이 집중도는 분포 진단일 뿐 의미 타당성 검증은 아니다.

\begin{figure}[H]
\centering
\begin{tikzpicture}
\begin{axis}[
  xbar, xmin=0, xmax=270, width=0.68\linewidth, height=11.2cm,
  xlabel={카드 수},
  symbolic y coords={Robot Control,Overconfidence,Weaponization,Non-Contestability,Conflict,Destabilising Dynamics,Excessive Authority,Network Effects,Privacy,Policy Exposure,Misinformation/Disinformation,Copyrights,Goal Misalignment,Anthropomorphism},
  ytick=data, nodes near coords, bar width=9pt, xmajorgrids, grid style={grid},
  tick label style={font=\small},
]
\addplot[fill=blue,draw=navy] coordinates {
 (45,Robot Control) (46,Overconfidence) (48,Weaponization) (55,Non-Contestability)
 (64,Conflict) (73,Destabilising Dynamics) (74,Excessive Authority) (78,Network Effects)
 (86,Privacy) (88,Policy Exposure) (107,Misinformation/Disinformation)
 (113,Copyrights) (150,Goal Misalignment) (242,Anthropomorphism)};
\end{axis}
\end{tikzpicture}
\caption{최종 정책 뷰의 상위 14개 L3. 정확한 조회를 위한 전체 50개 분포는 CSV 산출물에 제공한다.}
\end{figure}

\section{보류 55개와 불확실성}

\begin{longtable}{lrrr}
\toprule
보류 사유 & 건수 & 보류 내 비율 & 전체 비율 \\
\midrule
\endhead
PHYSICAL\_OUTSIDE\_LOCK & 23 & 41.8182\% & 1.3326\% \\
FRONTIER\_EXPERT\_REJECTED & 17 & 30.9091\% & 0.9849\% \\
FRONTIER\_EXPERT\_DISAGREEMENT & 7 & 12.7273\% & 0.4056\% \\
MULTI\_MECHANISM & 6 & 10.9091\% & 0.3476\% \\
LOW\_ABSOLUTE\_FIT & 2 & 3.6364\% & 0.1159\% \\
\midrule
합계 & 55 & 100.0000\% & 3.1866\% \\
\bottomrule
\end{longtable}

보류는 승인되지 않은 51번째 L3로 취급하지 않는다. 웹 UI와 export에서는 별도
검토 큐로 표시하고, 원래 보류 사유, Stage 3 강제 후보 L3, semantic/composite
점수, anchor vote를 숨기지 않아야 한다. 전문가 거절과 불일치 24개는 L3 정의의
적합성 검토, 다중 메커니즘 6개는 primary/secondary 또는 복수 태그 구조 검토,
Physical 경계 23개는 기존 182개 정답 집합과의 의미·근거 기반 비교가 필요하다.

보류 카드가 특정 영역에 집중되는 현상은 분류 오류만으로 설명되지 않는다.
다수의 L4 리스크 카드는 단일 상위 L3에 완전히 귀속되기보다 여러 위험 범주를
가로지르는 교차적 위험성(cross-cutting risk)을 보이며, 단일 부모 계층에서는
경계 사례(boundary case)로 나타난다. 이는 반드시 오분류라기보다 L3 내포 간
범주 중첩(category overlap)의 문제이다. 따라서 향후 확장에서는 필수 primary
L3와 선택적 secondary L3를 부여하는 다중 라벨 분류(multi-label
classification), 또는 하나의 L4가 복수의 L3 부모를 가질 수 있는 다중 부모
계층(polyhierarchy)으로 이들 카드를 수용하는 방안을 검토할 수 있다. 두 방안
모두 primary 부모를 필수로 유지하면 운영상 단일 책임 지점은 보존된다.

\section{검증, 제한과 강건성 점검}

\subsection{v2.17.2 역사적 기준선: L4 통폐합, 재정의, BGE-M3 재검증}

Claude의 최종 L4 명칭 정비, 일부 통폐합, 일부 재정의 결과를 v2.17.2에 반영했다.
ID 레지스트리는 1,711개 registered L4 ID로 불변이다. 이 중 1,660개는 active
card로 표시하고 검증하며, 51개는 \code{merged\_into} 이력을 가진 merged
provenance record로 보존한다. 따라서 사이트와 보고서는 두 수치를 항상
\textbf{1,711 registered IDs, 1,660 active cards} 형식으로 병기한다. 웹
탐색에서는 merged 배지를 통해 retired ID도 계속 확인할 수 있고, 신뢰성 검증은
active 1,660개를 기준으로 수행한다.

BGE-M3 기반 constrained-EM 매핑 알고리즘을 active 1,660개에 대해 다시 실행했다.
Physical 경로 189개는 잠금 대상으로 유지했다. 재실행 결과 144개의 guarded 이동
후보가 도출되었으며, 이들은 정답으로 확정하지 않고 모두 HOLD 검토 항목으로
반영했다. 기존 HOLD 카드는 \code{hold\_semantic\_path}만 갱신하고, 기존
비-HOLD 비Physical 카드는 General 또는 Agentic HOLD 검토 경로로 이동시킨 뒤
후보 semantic L3를 보존했다. 최종 active HOLD는 765개이다.

\begin{longtable}{lrrrr}
\toprule
지표 & 전체 사전 & 전체 사후 & 비HOLD 사전 & 비HOLD 사후 \\
\midrule
\endhead
평가 카드 수 & 1,660 & 1,660 & 954 & 954 \\
Top-1 일치율 & 65.5\% & 70.2\% & 77.3\% & 79.7\% \\
Top-5 일치율 & 92.9\% & 94.6\% & 96.0\% & 97.2\% \\
중앙 margin & 0.0116 & 0.0150 & 0.0258 & 0.0276 \\
음수 margin 비율 & 34.5\% & 29.8\% & 22.7\% & 20.3\% \\
수렴 스텝 & 15 & 15 & 19 & 19 \\
수렴 유사도 & 0.797 & 0.797 & 0.800 & 0.800 \\
노이즈 안정성, $\sigma=0.05$ & 72.0\% & 73.4\% & 78.8\% & 79.8\% \\
\bottomrule
\end{longtable}

결과적으로 L4 통폐합과 재정의 이후 전체 active 카드의 Top-1 일치율은
65.5\%에서 70.2\%로 개선되었고, 비HOLD 카드 기준 Top-1 일치율은 77.3\%에서
79.7\%로 개선되었다. 비HOLD 기준 노이즈 안정성은 78.8\%에서 79.8\%로
개선되었다. 다만 144개 guarded 이동 후보와 765개 active HOLD는 여전히 사람의
분류 검토가 필요한 항목이다.

Physical 182개와 Stage 1 합의 37개는 후속 단계에서 변경되지 않았다. Stage 2와
Stage 3 생성기는 동일 입력에서 byte-identical placement hash를 재생성했다.
Stage 2 validator는 13 PASS/0 FAIL/3 WARN, Stage 3 validator는
14 PASS/0 FAIL/4 WARN이었다. 최종 보류 정책 전용 테스트를 포함한 전체 회귀
테스트 16개가 통과했다.

주요 제한은 다음과 같다.
\begin{itemize}[leftmargin=2em]
  \item Stage 2 신규 1,334개 중 1,214개는 rule eligibility 없이 의미 top-1 중심으로 배치됐다.
  \item Stage 2 신규 247개는 taxonomy-gap sentinel을 override했다.
  \item Stage 3에 남긴 118개도 calibrated probability나 사람 승인 결과가 아니다.
  \item Physical 경계 탐지는 제한된 키워드 sentinel이므로 오탐과 미탐 가능성이 있다.
  \item 50개 L3의 사용률 100\%와 카드 커버리지 96.8134\%는 분류체계의 충분성을 입증하지 않는다.
\end{itemize}

\subsection{v2.18.0-rc 계층 확장과 HOLD 민감도 검증}

현재 릴리스 후보는 General AI에 사회적 파급 L2를 추가하고, 에너지 소비 및
환경 오염, 책임성 부족 및 거버넌스 체계 부재, 투명성 부족, 공정성의 4개 L3를
추가했다. 기존 Physical AI의 사회적 안정성 관련 L3는 General AI 사회적 파급
아래로 이동했다. 최종 의미 L3는 General 33개, Agentic 6개, Physical 15개로
총 54개이다. 1,711개 등록 ID와 1,660개 활성 카드, 51개 통합 이력은 보존했다.

동일한 BGE-M3 임베딩과 54개 L3 중심점을 사용해 HOLD 포함 1,660개와 HOLD 제외
1,046개를 비교했다. HOLD 제외 시 수렴 스텝은 20회에서 12회로 감소했고,
평균 코사인 목적함수는 0.8005에서 0.8026으로 증가했다. Top-1 포함률은
71.9\%에서 81.6\%로 9.7\%p, $\sigma=0.05$ 노이즈 안정성은 75.0\%에서
80.7\%로 5.7\%p 개선됐다.

\begin{longtable}{lrr}
\toprule
검증 지표 & HOLD 포함 & HOLD 제외 \\
\midrule
\endhead
평가 카드 수 & 1,660 & 1,046 \\
의미 L3 수 & 54 & 54 \\
수렴 스텝 & 20 & 12 \\
평균 코사인 유사도 & 0.8005 & 0.8026 \\
Top-1 포함률 & 71.9\% & 81.6\% \\
Top-3 포함률 & 91.3\% & 94.6\% \\
Top-5 포함률 & 95.2\% & 97.2\% \\
노이즈 안정성, $\sigma=0.05$ & 75.0\% & 80.7\% \\
\bottomrule
\end{longtable}

\begin{figure}[H]
\centering
\includegraphics[width=\textwidth]{../validation/v2.18.0-rc/hold_sensitivity_bge_m3/hold_included_excluded_sensitivity_3panel.pdf}
\caption{HOLD 포함 여부에 따른 민감도 검증. 이미지 제목에는 릴리스 번호를
표시하지 않고, 보고서 캡션과 기계 판독 메타데이터에서 버전을 관리한다.}
\end{figure}

\section{권고 검토 순서와 후속 질문}

\begin{enumerate}[leftmargin=2em]
  \item 현재 HOLD 614개를 General과 Agentic 영역, HOLD 원인, 의미 적합도 구간에 따라 층화하여 검토한다.
  \item 신규 또는 이동된 L3의 경계 사례를 우선 표집해 휴먼 평가로 Top-1 배치의 타당성을 검증한다.
  \item General 사회적 파급으로 이동한 9개 L3와 신규 4개 L3의 정의 중복 및 배타성을 검토한다.
  \item 다중 메커니즘 카드에는 primary L3를 유지하면서 secondary L3 또는 다중 부모 구조의 필요성을 평가한다.
  \item 사람 승인을 받은 변경만 후속 릴리스와 홈페이지에 반영하고 1,711개 ID 이력을 보존한다.
\end{enumerate}

추가로 결정해야 할 질문은 (1) HOLD 해제에 필요한 사람 합의 기준, (2) 다중
메커니즘 카드의 secondary L3 허용 범위, (3) 사회적 파급 L3 간 경계 정의,
(4) 후속 신규 L3 제안의 최소 증거 기준이다.

\vfill
\noindent\textbf{보고서 식별자:} RAI Risk Taxonomy Technical Report 2.0\\
\textbf{정책 데이터 뷰:} \code{v2.18.0-rc}\\
\textbf{작성일:} 2026-07-25

\end{document}
